\subsection{Мощность теплопередачи}

%%%%%%%%%%%%%%%%%%%%%%%%%%%%%%%%%%%%%%%%%%%%%%%%%%%%
%%%%%%%%%%%%%%%%%%%%%%%%%%%%%%%%%%%%%%%%%%%%%%%%%%%%
 \z{В сосуд с горячей водой массой $m = 0.5\,\text{кг}$ опустили работающий нагреватель. В результате температура воды повысилась на $\Delta T = 1\,\degree\text{C}$ за время $t_1 = 100\,\text{с}$. Если бы воду не нагревали, то её температура понизилась бы на ту же величину $\Delta T$ за время $t_2 = 200\,\text{с}$. Какова мощность нагревателя?}
%
{$P = c m \Delta T \left(\frac{1}{t_1} + \frac{1}{t_2}\right) = 31.5\,\text{Вт}$}

%%%%%%%%%%%%%%%%%%%%%%%%%%%%%%%%%%%%%%%%%%%%%%%%%%%%
 \z{Школьник Владислав исследует охлаждение воды в стакане на морозе. Владислав заметил, что охлаждение от температуры $91\,\degree\text{C}$ до $89\,\degree\text{C}$ происходит за $3\,\text{минуты}$, а от температуры $31\,\degree\text{C}$ до $29\,\degree\text{C}$ -- за $6\,\text{минут}$. Чему равна температура окружающей среды? Считайте, что мощность теплоотдачи пропорциональна разности температур стакана и окружающей среды.}
%
{$-30\,\degree\text{C}$}

%%%%%%%%%%%%%%%%%%%%%%%%%%%%%%%%%%%%%%%%%%%%%%%%%%%%
 \z{На краю крыши висят сосульки конической формы, геометрически подобные друг другу, но разной длины. После резкого потепления от $T_1 = 0\,\degree\text{C}$ до $T_2 = 10\,\degree\text{C}$ самая маленькая сосулька длиной $l = 10\,\text{см}$ растаяла за время $t = 2\,\text{часа}$. За какое время растает большая сосулька длиной $L = 30\,\text{см}$, если внешние условия не изменятся?}
%
{$6\,\text{часов}$}

%%%%%%%%%%%%%%%%%%%%%%%%%%%%%%%%%%%%%%%%%%%%%%%%%%%%
 \z{При нагревании на плите кастрюли с водой температура воды увеличилась от $90\,\degree\text{C}$ до $92\,\degree\text{C}$ за 1 мин. Какая доля энергии, получаемой водой, рассеивается в окружающем пространстве, если время остывания той же воды от $92\,\degree\text{C}$ до $90\,\degree\text{C}$ равно 9 мин?}
%
{$\alpha = 0.1$}

%%%%%%%%%%%%%%%%%%%%%%%%%%%%%%%%%%%%%%%%%%%%%%%%%%%%
 \z{ Если температура на улице равна $t_1 = -20\,\degree\text{C}$, то температура в комнате устанавливается равной $t_2 = +20\,\degree\text{C}$. Если же на улице похолодает до $t_3 = -40\,\degree\text{C}$, температура в комнате станет равна $t_4 = +10\,\degree\text{C}$. Найдите температуру комнатной батареи.}
%
{$t= 60 \,\degree\text{C}$}

%%%%%%%%%%%%%%%%%%%%%%%%%%%%%%%%%%%%%%%%%%%%%%%%%%%%
 \z{Горячий суп, налитый доверху в большую тарелку, охлаждается до температуры, при которой его можно есть без риска обжечься, за время $t = 20\,\text{мин}$. Через какое время можно будет есть суп с той же начальной температурой, если разлить его по маленьким тарелкам, которые также заполнены доверху и подобны большой? Известно, что суп из большой тарелки помещается в $n = 8$ маленьких, и что количество тепла, отдаваемое в единицу времени с единицы поверхности каждой тарелки, пропорционально разности температур супа и окружающей среды.}
%
{$10\,\text{минут}$}

%%%%%%%%%%%%%%%%%%%%%%%%%%%%%%%%%%%%%%%%%%%%%%%%%%%%	
\z[{\begin{wrapfigure}[8]{r}{0.3\textwidth}%
		\vspace{-15pt}%
		\begin{flushright}%
 			\includegraphics[width=0.3\textwidth]{1.5.5}%
		\end{flushright}%
	\end{wrapfigure}%
	}]%
	{Система охлаждения нагревателя состоит из нескольких одинаковых теплопроводящих стержней, соединённых небольшими шариками. Температура нагревателя $70\,\degree\text{C}$, температура холодильника $35\,\degree\text{C}$. Чему равна разность температур шарика $A$ и шарика $B$ в установившемся режиме? Теплопроводящая система теплоизолирована. Приток тепла осуществляется только от нагревателя, а отвод только через холодильник. Считать, что мощность теплопередачи через стержень пропорциональна разности температур на его концах.}
%
{$15\,\degree\text{C}$}



